% Independent mathematical review; source originals are unchanged.
\section{Independent review of the strong C2, IH, HI and UI closure}
\label{sec:independent-strong-c2-review}

\paragraph{Scope and conclusion.}
This review reads the mathematical bodies of components
40--43, 64, 67--78, 83--87, 96 and the corresponding joins in 194.
It does not treat a historical verification receipt as evidence.
Within the stated physical native C2 source interface, no mathematical
defect was found in these retained branches. The conclusion is about
their complete tensor-SOS constructions, including all state normals,
source openings, arbitrary independent testing rates and literal
$n=0$. It is not a claim that their scalar graph identities alone are
sufficient. Native canonical entry, the detailed noncommuting-testing
theorem A6, matching/two-weak entry, the single-angle chain, the weak-II
HJ/BSR chain, and the positive-total-scale terminal are separate inputs
reviewed by other members of this audit.

The fresh script \texttt{strong\_c2\_independent.py} constructs the
source formulas and the homogeneous parent independently of the archive
checkers. Its 53 successful exact checks and denominator residues are recorded in the
adjacent JSON file. Those checks supplement the order, boundedness,
cap and coverage arguments below; they do not replace them.

\subsection{Common conventions and source estimates}

Let $K_c,K_t$ be independently quantified real closed extensions of
$\mathbb R$. A source division or radical belongs only to $K_c$;
a testing radical belongs only to $K_t$. A certificate row may be
unbounded and may depend on the other rows. Only the final Gram
coefficients must belong to the bounded tensor ring and have a strictly
positive definite real residue. The foundational strict-Gram lemma then
gives an actual finite tensor SOS. Finite real caps are chosen after the
source and testing residues have been fixed. These quantifiers are
compatible with Corollary 3.19 and supply no uniform block-size bound.

Write $z=u_c^2$, $y=v_c^2$, $H=zh$, $L=yt$, $B=H+L$ and
$A_c=1+zr$. Component72 gives the exact six source coefficients GE1.
Its positive decomposition
\[
 M=(1+z)(B+zr)H+y\{B+z(H+r-1)\}+yzj+zrL(1+y)
\]
is especially useful: it does not estimate a difference of positive
terms. For infinitesimal $z$, strong $C,\alpha$ force $B$ to be a
bounded positive unit. The inequalities $CL=O(1)$ and
$\alpha(1-H)\le(B+zr)(1+z)/2$ prove the upper and lower assertions.
Nonnegativity of $n$ gives $yr\le A_c(A_c+L)$, hence the denominators
used in GE2 are units after their displayed compensating factors.
The exact identity
\[
 [(B+zr)A_c(A_c+L)-r\delta]_{r=1}
 =-(H+y)(B+z)(1+z)
\]
proves that $1-r$ is a unit; this argument requires neither a lower
bound on $r$ nor strict positivity of $n$.

For weak-$C$ IH, component77's WI1--WI3 show instead that $B$ is
unbounded. If $B$ were bounded, $\widehat N/[y(B+zr)]\to1$ and
$\widetilde\Delta/[y(B+zr)]$ would be bounded, forcing $C\ge c/L$
for a positive real $c$. With $B$ unbounded, $\rho=B/y$ has residue in
$[0,1]$, $r/\rho$ is bounded, and $1-r$ is a unit. The five normalized
products in WI3 absorb every occurrence of $H$, $1-H$ and the endpoint
gap $e=1+zr-L/y$. The illustrative curves and former frontier statements
in that component are not premises of the closure.

\subsection{Strong II: the dependent row and the missing angular pivot}

The local construction in component40 uses the graph
\[
 X=-zr vW/B+\xi,\qquad Y=HuW/B+\eta,\qquad Z=uvW/B+\zeta.
\]
This is a polynomial change of all four homogeneous state coordinates.
It differs from the graph in component64; the latter graph is not
needed by this local proof. The harmonic completion has denominator
$cj+dt$, which is a positive unit because $c,d$ are units and
$t+j=A_c$ is a unit. Thus the opening limits $t=0$ or $j=0$ in residue
do not insert a singular bounded-Gram coefficient.

The delicate step is correctly handled by LII7. The row
$R_h=(vQ_n-su(V-y)W)/k$ is dependent. Before using the physical state
jet, the entire formal $R_hW$ coefficient is eliminated using
\[
 (ukW)R_h=uvWQ_n-sU(V-y)W^2.
\]
The $\eta$ derivative then forces the formal $N$ coefficient to vanish,
and the $\xi,\zeta$ derivatives force the two remaining formal
coefficients. Without this gauge, physical stationarity would not by
itself imply all the separate formal coefficient vanishings claimed in
LII9. With the gauge, LII9--LII10 are valid.

The support statement $P_\xi\in(U,z)$ is bounded-coefficient ideal
membership: its possible source factors are $m=O(z^2)$,
$nzr=z(nr)$ and $zr$. Monic division therefore produces
$P_\xi=(U-z)A_\xi+z(V-y)B_\xi$ without a hidden small inverse.
The corresponding scalar decomposition LII13 follows by two monic
divisions and the exact value and angular first derivatives. LII15
then allocates every scalar and first-normal term on the six rows
$(N,\xi,\zeta,R_h,R,T)$. No state normal or higher testing term is
discarded.

The five-row block excluding $T$ is strict: its only nontrivial
two-row Schur complement is $2/[h(1+h)^2]>0$. Finally the polynomial
zero LII20 adds arbitrary positive real mass to $T^2$ while changing
all other residue entries by zero. A sufficiently large real multiple
therefore closes the final Schur complement. This is lawful even when
$r_0=0$, $n=0$ or $h=1$.

Component41 treats the original homogeneous parent directly. It does
not homogenize the higher-degree local graph. Its remote harmonic
frames have only bounded coefficients; for the remaining two blocks,
RR10 is a nonempty real interval because $0<b_0<1$ and $h_0$ can be
chosen sufficiently close to $h$. At $b_0=0$ the exact three-square
identity RR13 replaces that interval argument. Thus literal projective
endpoints, arbitrary states and all source openings are covered.

\subsection{IH: one exact current calculation and three source shapes}

Components76 and75 retain the entire reciprocal-Bob parent IC3.
The graph $X=-ruZ$, $Y=HwZ$ has two unrestricted state normals.
The five currents in IC5, the exact scalar identity IC7, both toric
identities IC8/IC8a, and the full normal formulas IC10--IC12 are
polynomial identities at the original source. Their use by the strong
IH branch imports these identities only, not the weak-source bounds.

An especially helpful cancellation makes the endpoint estimates
transparent. Put $q=H+r-1$ and $S=A_c+B-H$. With $M,\delta$ as in GE1,
the coefficient $k$ in IC6 is exactly
\[
\begin{split}
k={2\over\delta}\bigl[&HA_c(B+zr)S(y+2)
 +(1-H)(B+zr)M\\
 &+(1+y)(1-r)A_cS(zq-B)\\
 &+r(1+2z)\{(B+zr)A_cS-r\delta\}+q^2\delta\bigr]
 +8DHr.
\end{split}
\]
Indeed the two $m$ terms combine by
$z^2q^2-B^2=(zq-B)(zq+B)$ and $zq+B=G$. This cancels the apparent
$1/G$ before any source endpoint is taken.

For strong IH, substitute $y=B/\rho$, $r=s\rho$ with $B$ a unit.
For weak IH at $\rho_0=0$, substitute
$B=1/\tau$, $y=1/(\tau\rho)$, $r=s\rho$.
For weak IH at $\rho_0=1$, use
\[
 B=1/\tau,\qquad
 \rho={1+zr-e\over1-H\tau},\qquad y={1\over\tau\rho}.
\]
The independent calculation cancels the resulting rational functions
before specializing the simultaneous infinitesimals. Every displayed
normalized source quantity in GE3, WIE1 and WIE3 has a nonzero
denominator residue, including $H_0=0,1$. In particular,
\[
 \left({k\over B\rho}\right)_0
 ={2(B^2s+BHs+BH-Bs-H^2+H)\over B^2}
\]
in strong IH, it is $2s$ in the weak $\rho=0$ case, and
$((k-8DHr)/B)_0=r$ at the weak $\rho=1$ endpoint.
These are simultaneous limits, not iterated limits. They exclude the
possible hidden ratios $z/\rho$ and uncompensated opening poles.

For interior weak IH, the positive term added by IC8a gives the $T$
row a unit pivot even at $r_0=0$. The displayed two-row determinant
is positive for $0\le r\le\rho/(1+\rho)$; the other scalar pivots
are positive. The full normals are divided by
$\sqrt{4n+2}$ and $\sqrt{B/[H(1-H)]}$. The exact quotient formulas
IC11--IC12 make all their scalar couplings infinitesimal, while the
normal residue is $\operatorname{diag}(1,2/(1+\rho))$. This uses the
compensated products $AH/B$ and $\alpha(1-H)/B$, so both opening
endpoints are included.

At $\rho_0=0$, replace $T$ by $T_\rho=\rho T$ and multiply IC8a
by $\rho$. The remaining two-row block is
\[
 \begin{pmatrix}2&-s\\-s&1+2s\end{pmatrix},\qquad 0\le s\le1,
\]
whose determinant is at least two. The coefficient $f_X/\rho$
remains bounded, so this rescaling does not lose the first $X$ normal.
At $\rho_0=1$, retain the whole full-state square
$4D(J+eZ)^2$. The D-removed normal quotients have only bounded
$De$ constants. The exact relation WIE5 eliminates $R_2$ using the
pure-source infinitesimal $1/\sqrt D$. The final seven-row residue
is the strictly positive diagonal in WIE6. Subtracting only a scalar
graph square instead would not justify this step; the retained square
here is explicitly a full-state square.

For strong IH, the entire explicit $C$ dependence is precisely the
two full-state squares SL2. Retaining their excess and capping $C$
does not alter the original coefficient $k$. The two finite blocks
have determinants $8c_*(1+L)$ and $2a_*c_s$, both positive at all
allowed openings. A real multiple of the scaled IC8a identity buys
the remaining $T_\rho$ pivot. This establishes the whole parent,
not only the graph.

\paragraph{Remote IH and the open axis.}
Component74's remote corners include the entire normal restoration;
their $\sqrt d$ rows never create an inverse of a weak $d$ in a
Gram coefficient. Component69 is indispensable. On the open axis
$u_0\ne0$, $w\sim0$, its exact identity $PE=VZ$ splits the cross
term between two blocks. The interval OA7 is nonempty for every
$s=\sqrt{1+u_0^2}>1$ and $0<\rho_0<1$. At $\rho_0=1$ a finite
real cap on the unbounded $D$ is sufficient; at $\rho_0=0$ the entire
$D$ packet is retained and the unscaled $Z$ row disappears from the
exact remainder. In particular its first $X$ coefficient then is
$uV+u(1-H)wE/\sqrt B$, with no omitted $Z$ dependence.

The strong remote theorem71 treats the same endpoint forms with
its own source estimates. At opening boundaries its real penalty
argument is valid because the constrained residue forms are explicitly
strict. In the both-infinite testing chart the core contribution to
the $E^2$ pivot is retained; the potentially finite normal coefficient
$2\{AH-\alpha(1-H)\}-(1-H)$ has zero residue. Its sole external
testing input is A6 in the unit/unit cell. Together these bodies cover
the full nine testing-type cells; the stand-alone join73 adds no new
mathematical ingredient.

\subsection{HI: source exclusion, four-term circuit and arbitrary state}

Components87 and67 derive the source normalization from $n\ge0$.
Their positive remainder in WE3 makes the contradiction for unbounded
$R=r/q$ rigorous: dividing by $r^2(H+yr)$ leaves a positive unit
minus infinitesimals. If $r$ is weak, the gap $g>0$ forces $H_0=1$;
if $r$ is a unit, the necessary quotient is bounded by $1/r$.
This handles both alternatives and literal $n=0$.
Consequently $R$ is bounded, $H\sim1$, and all four coefficients
$m,C,D,\alpha$ are unbounded. Component68 can therefore choose
arbitrarily large finite real caps after fixing the testing residue.
Its strict kernel/penalty forms cover the entire complement of the
critical nonzero infinitesimal $p=1/u,v$ corner, including its literal
axes.

At the critical corner, component43's state coordinates
$(s,Q,N_c,N_m)$ are a full linear state parametrization with only
pure-source divisions. Both $s$ and $Q$ are decomposable tensors.
The sign of $sQ$ is determined by a pure-testing expression $-SY$,
so the case split and every positive fourth root in the circuit are
lawful. No order or radical is taken on a general mixed tensor.
The graph roots $A=Ps$, $B=Vs$, $C=PVQ$, $D=Q$ and
$G=(ABCD)^{1/4}$ satisfy the exact finite circuit decompositions
N0M8. A finite positive average supplies simultaneous strict pivots
on every listed variance and conductor row.

The exact half-angle identities
\[
 G(\sqrt p-q^{1/4})=q^{1/4}(\sqrt{AC}-G),\qquad
 G(\sqrt v-y^{1/4})=y^{1/4}(\sqrt{BC}-G)
\]
are the decisive conductors. Finite polynomial Taylor expansion in
$\sqrt p,\sqrt v$ has bounded coefficients. Its degree-two and
higher terms become infinitesimal bounded quadratic forms on the
conductor rows. The potentially dangerous degree-zero and degree-one
$G$ terms cancel by the exact state and angle jets N0M6. The displayed
identity for that cancellation is algebraically correct. Both first
state-normal sums vanish at contact and are allocated by the same
conductors; their complete raw terms are present.

For arbitrary $n\ge0$, component42 retains the entire actual packet
$4n d_p d_v x^2$, with its centered full-state root $x$.
The remaining coefficients depend on $n$ only through the bounded
$\nu=2nR$. All four graph coefficients have common residue
$2(1+r)^2/(3+r)\in[2/3,2]$, and both normal pivots are
$4/(3+r)>0$, including $r=0$. Thus the predecessor's finite circuit
argument applies without an inference from $n=0$ to $n>0$.
For $sQ<0$ or $sQ=0$ the raw-root diagonal Gram is already strict.
The proof is a pointwise tensor SOS for each testing state; the roots
need not remain linear polynomials in state coordinates. That is
exactly sufficient for the qualitative criterion being used.

\subsection{UI: all normals, the matched Hessian and finite caps}

Component83 supplies the entire centered-packet identity UR16.
In particular, it cancels $Ct-Dj$ before estimating the harmonic
normal and keeps every signed-core term. The denominators of the
compensated products in UR9--UR11 are positive units. Their continuity
through allowed residue openings is therefore algebraic.

The local proof84 gives a polynomial inverse for all state normals.
Its scalar and first-normal divisions are monic in $u-a$ and
$v^2-k^2$. The identities UI1 allocate them without a mixed inverse.
The $E_r$ first-normal coefficients have zero residue. The remaining
bounded Gram has six rows, with the ordinary-angle rows $S,T$
separate. Setting $S=T$ gives the matched full leading Hessian after
an invertible bounded coordinate change. The Schur formulas UI4--UI5
are checked directly from the independently rebuilt RC4 leading
polynomial, including both state directions before elimination.
Their positivity follows from
\[
 Q_n=(1-(z+1)r)(1+zr)-zr g\ge0
 \quad\Longrightarrow\quad r\le1/(z+1).
\]
Every denominator in these Schur formulas is a unit. The exact zero
in UI1 subsequently distributes the positive total Schur mass between
$S^2$ and $T^2$; its other residue contributions vanish.

For diverging $n,m,\alpha$, the cap argument is sound but requires
the real-interior approximation supplied in component78. It is not
lawful to substitute an actual $b$ directly into a leading formula
with denominator $1-zb$ when that denominator changes sign at an
infinitesimal rate. Instead use
$b_s=sb+(1-s)/(1+z)$, $r_s=sr$ and the exact identity
\[
 f(b_s,r_s)=(1-s)(1+s-sr)+s^2f(b,r).
\]
Small further perturbations make the remaining gaps strict. This
approaches every allowed boundary by valid finite real leading
sources. At most two hard-state penalties diverge, and every possible
pair of their hard rows is independent with a unit determinant.
One can therefore adjust the hard coordinates into their exact
nearby kernels. The positive matched Schur form and bounded minimizing
hard vector give a uniform positive lower bound on the complete
Hessian. Taking the kernel limit and applying the real finite-penalty
lemma yields lawful finite caps below the original coefficients.

The ordinary-angle complement in78 is genuinely all-state. The
exact relation $R_0K_0-(uX-Z)V_0=0$ permits the real rank-one
allocation UA1--UA3. The normalized Laurent determinant UA11 has
nonnegative nonconstant coefficients, including strictly positive
$U^{-1}$ and $T$ coefficients. In logarithmic coordinates these
already span both directions, so its stationary zero is the unique
minimum and it is strictly positive away from the contact. UA12
shows that the inverse expressions used only in this real residue
calculation extend through the divergent source boundaries. The
same valid real-interior approximation justifies the cap limit there.

Every remaining UI endpoint is treated by the original homogeneous
parent in UA13--UA16. The displayed forms have strict residue in
their stated complete roots; roots themselves may be unbounded.
Literal axes and projective endpoints require no testing division.
Together with the local certificate and the ordinary complement this
closes UI and, by exact reflection, OU.

\subsection{Angle exclusions, five-event terminal and assembly}

Component96's ordinary/ordinary argument shows $m,A,D$ strong and
proves that weak $C$ or weak $\alpha$ forces both complementary
events weak. Its three directed exclusions use exact source
inequalities and positive numerators; the small/ordinary exclusion
also uses the separately audited bounded-state/unbounded-$b$ entry.
Component87 rules out weak $C$ at an unbounded first actual angle
and a bounded second one. These exclusions concern actual source
angles only and impose no testing-rate restriction.

In component86, the selected five-event capped parent is used only
on commuting testing strata. NW1--NW5 give strict full-state or
coordinate Hessian forms there, and every higher monomial has an
infinitesimal bounded Gram allocation. The noncommuting residue case
uses A6 on the original physical source; the capped parent is never
assumed physical. This is an essential distinction and is respected
by the proof.

The source table and reflection in194 assemble these branches with
the separate matching, single-angle and weak-II chains. All source
types and independent testing types are exhausted. No branch above
uses the global C2 conclusion as a premise. Its old frontier remarks
are historical scope descriptions, not mathematical hypotheses of
the retained local identities.

\paragraph{Minimal closure and honest compression.}
The exact equation-level closure and safe deletions are in
\texttt{strong\_c2\_minimal\_closure.json}. In particular, 64's
unused alternate graph, 73's duplicate join, example curves, old
frontier paragraphs and verification-status prose can be removed.
The all-$n$ HI proof should merge43 and42. The UI boundary
approximation in78, the IH open-axis proof69, the full normal
formulas and the source endpoint calculations cannot be removed.
This review fragment is a mathematical audit and compression map.
It is not a replacement for the exact certificate equations named
above unless those equations are incorporated into the refined
manuscript's explicit certificate appendix.
